% Cálculo del tensor de estrés en unav-walkers: simulación y procesamiento.
% Compilar con: latexmk -pdf estres.tex
\documentclass[11pt,a4paper]{article}

\usepackage[T1]{fontenc}
\usepackage[spanish,es-noshorthands,es-tabla]{babel}
\usepackage{lmodern}
\usepackage[margin=2.5cm]{geometry}
\usepackage{amsmath,amssymb}
\usepackage{booktabs}
\usepackage{array}
\usepackage{tikz}
\usepackage[hidelinks]{hyperref}

\newcommand{\vect}[1]{\mathbf{#1}}
\newcommand{\n}{\vect{n}}
\newcommand{\tv}{\vect{t}}
\newcommand{\lv}{\vect{l}}
\newcommand{\f}{\vect{f}}
\newcommand{\vv}{\vect{v}}
\newcommand{\code}[1]{\texttt{#1}}
\newcommand{\ave}[1]{\langle #1 \rangle}

\title{Tensor de estrés en \code{unav-walkers}:\\
  cálculo en la simulación y procesamiento posterior}
\author{Manuel Carlevaro}
\date{Versión del código: simulador 3.3, \code{stress\_profile} con errores
  de la parte cinética}

\begin{document}
\maketitle

\begin{abstract}
  Este documento describe todo el cálculo relacionado con el tensor de
  estrés: cómo se reconstruyen las fuerzas de contacto a partir de los
  impulsos de Box2D, cómo el simulador arma el tensor de cada grano y lo
  escribe en los archivos \code{.sxy}, cómo \code{stress\_profile} los
  promedia en una franja vertical para obtener perfiles en $y$ (parte de
  contacto, normal, tangencial y cinética, con sus errores) y cómo
  \code{plot\_stress\_profile.py} los convierte a unidades experimentales.
  Cada fórmula remite a la función del código que la implementa.
\end{abstract}

\tableofcontents

% ---------------------------------------------------------------------------
\section{Flujo de datos}

\begin{center}
\small
\begin{tabular}{>{\raggedright}p{0.22\textwidth}>{\raggedright}p{0.3\textwidth}p{0.36\textwidth}}
  \toprule
  Etapa & Código & Resultado \\
  \midrule
  Fuerzas de contacto & \code{contact\_point\_force} (\code{src/contacts.cpp})
      & $F_n$, $F_t$, $\n$, $\tv$ en cada punto de contacto \\
  Tensor por grano & \code{save\_stress} (\code{src/output.cpp})
      & \code{<pre>\_<frame>.sxy}: $s$, $s^n$, posición, velocidad \\
  Perfiles en $y$ & \code{tools/stress\_profile.cpp}
      & $s$, $s^n$, $s^t$, $k$, $\phi$, errores por bloques \\
  Gráficos & \code{plot\_stress\_profile.py} y \code{unidades.py}
      (en \code{tools/scripts/}) & seis paneles en unidades experimentales \\
  \bottomrule
\end{tabular}
\end{center}

% ---------------------------------------------------------------------------
\section{Convenciones}

\paragraph{Unidades.} La simulación usa unidades reducidas: el diámetro
$d$ de los discos, su masa $m$ y el tiempo $\sqrt{d/g}$ valen 1. La
fuerza se mide en $mg$ y el estrés bidimensional (fuerza por unidad de
longitud) en $mg/d$. La conversión a las unidades del experimento está en
la sección~\ref{sec:unidades}.

\paragraph{Signo.} El estrés es \emph{negativo en compresión}. En un
empaquetamiento comprimido, $s_{xx}$ y $s_{yy}$ son negativos; por eso se
grafican $-s_{xx}$, $-s_{yy}$ y la presión $p = -(s_{xx}+s_{yy})/2$.

\paragraph{Orden de los índices.} En todo el código,
\begin{equation}
  s_{ij} \propto \sum f_i \, l_j ,
\end{equation}
con el primer índice en la fuerza $\f$ y el segundo en el brazo $\lv$. El
orden solo importa para la parte antisimétrica (sección~\ref{sec:antisim}). Las
componentes se escriben en el orden \code{xx xy yx yy}.

\paragraph{Ejes.} El silo ocupa $|x| \le R$, $0 \le y \le H$, con el
orificio en $y = 0$, $|x| \le r$. El eje $y$ de la simulación es opuesto
al de los experimentos.

% ---------------------------------------------------------------------------
\section{Fuerzas de contacto a partir de Box2D}
\label{sec:fuerzas}

Box2D (versión 2.4.2) resuelve los contactos con impulsos. En cada paso
de tiempo $\Delta t$, el solver de velocidades acumula en cada punto del
manifold un impulso normal $\lambda_n \ge 0$ y uno tangencial
$\lambda_t$, con $|\lambda_t| \le \mu \lambda_n$. La fuerza media durante
el paso es
\begin{equation}
  F_n = \frac{\lambda_n}{\Delta t}, \qquad
  F_t = \frac{\lambda_t}{\Delta t}.
\end{equation}

\paragraph{Normal y tangente.} La normal $\n$ del contacto apunta del
cuerpo A al cuerpo B. La tangente es la que usa \code{b2ContactSolver}:
\begin{equation}
  \tv = \code{b2Cross}(\n, 1) = (n_y, -n_x).
\end{equation}
La fuerza sobre B es
\begin{equation}
  \f_B = F_n \, \n + F_t \, \tv , \qquad \f_A = -\f_B .
  \label{eq:fuerza}
\end{equation}
Esta convención está centralizada en \code{contact\_point\_force}; todas
las salidas y los diagnósticos la usan. Se validó con el balance de
impulso de cada grano entre dos pasos: el residuo relativo es $\sim
10^{-5}$ (lineal) y $\sim 10^{-6}$ (angular), contra $\sim 0.1$ con la
tangente invertida.

\paragraph{Punto de contacto.} Para dos discos, Box2D ubica el punto de
contacto en el punto medio entre las superficies, sobre la línea de los
centros. Si $\delta$ es la superposición, el brazo desde el centro de cada
disco es
\begin{equation}
  \lv_A = +\rho_A \, \n, \qquad \lv_B = -\rho_B \, \n, \qquad
  \rho = R - \delta/2 .
  \label{eq:brazo}
\end{equation}
Lo mismo vale para un disco contra una pared (cadena o segmento): el punto
está sobre la normal que pasa por el centro del disco. En los dos casos,
$\lv \parallel \n$. Esta propiedad, exacta para discos, es la que da las
identidades de la sección~\ref{sec:propiedades}.

\begin{center}
\begin{tikzpicture}[scale=1.6,>=stealth]
  \draw[thick] (0,0) circle (1);
  \draw[thick] (1.9,0) circle (1);
  \fill (0,0) circle (0.03) node[left] {A};
  \fill (1.9,0) circle (0.03) node[right] {B};
  \draw[->,gray] (0,0) -- (0.93,0);
  \node[gray,below] at (0.45,0) {$\lv_A$};
  \draw[->,gray] (1.9,0) -- (0.97,0);
  \node[gray,below] at (1.45,0) {$\lv_B$};
  \fill (0.95,0) circle (0.035);
  \draw[->,very thick,blue] (0.95,0.07) -- (1.6,0.07) node[above] {$\n$};
  \draw[->,very thick,red] (0.95,0) -- (0.95,-0.65) node[below] {$\tv$};
  \node[align=left,font=\small] at (4.6,0)
    {$\n$: de A hacia B\\ $\tv = (n_y, -n_x)$\\ $\f_B = F_n\n + F_t\tv$};
\end{tikzpicture}
\end{center}

\paragraph{Qué no incluyen las fuerzas.}
\begin{itemize}
  \item \emph{Impulsos de tiempo de impacto (TOI).} Box2D no los guarda en
    los manifolds. Por eso \code{continuous\_physics} está desactivado por
    defecto: con TOI activo, las fuerzas registradas quedarían incompletas.
  \item \emph{Corrección de posición.} Las iteraciones de posición mueven
    los cuerpos para reducir la superposición sin transferir impulso; no
    aportan a $\lambda_n$ ni a $\lambda_t$.
\end{itemize}

\paragraph{Coeficientes.} Box2D combina la fricción de dos cuerpos como
$\mu = \sqrt{\mu_1 \mu_2}$ y la restitución como $\max(e_1, e_2)$.

\paragraph{Instante de la medición.} En cada iteración del lazo principal
(\code{src/main.cpp}) primero se aplica la fricción de la base y se
escriben las salidas, y después se llama a \code{Step}. Un archivo con
tiempo $t$ contiene las posiciones y velocidades al final del paso
anterior y los impulsos de ese paso, es decir, las fuerzas medias en
$[t - \Delta t, t]$.

% ---------------------------------------------------------------------------
\section{La base vibrada}
\label{sec:base}

Los discos se apoyan sobre una base horizontal que vibra en la dirección
$y$ del silo. La gravedad no actúa en el plano: solo fija la carga normal
$N = mg$ sobre la base. La aceleración de la base es
\begin{equation}
  a(t) = \Gamma g \left[ \rho \sin(\omega t) + (1-\rho) \sin(2 \omega t +
  \varphi) \right], \qquad \omega = 2 \pi f ,
\end{equation}
y la base se mueve con $-y(t)$ en coordenadas del silo. La fase que se
registra en la cabecera de cada archivo es $\omega t \bmod 2\pi$.

La base ejerce sobre cada grano (\code{src/base\_friction.cpp}):
\begin{itemize}
  \item una fuerza de fricción de Karnopp, aplicada en el centro, que
    depende de la velocidad relativa a la base, con adherencia hasta
    $\mu_s N$ y deslizamiento con $\mu_d N$;
  \item un torque de pivoteo de Coulomb que se opone al giro, con límite
    estático $\frac{2}{3} \mu_s N R$ y valor dinámico
    $\tau_\text{piv} = -\frac{2}{3} \mu_d N R \, \operatorname{sign}\omega$.
\end{itemize}
Estas fuerzas \emph{no} entran en el tensor de estrés, que se calcula solo
con los contactos. El torque de pivoteo, sin embargo, deja una huella en
su parte antisimétrica (sección~\ref{sec:antisim}).

% ---------------------------------------------------------------------------
\section{Tensor de estrés de cada grano}
\label{sec:grano}

\subsection{Definición}

Para cada grano $p$, \code{save\_stress} suma sobre todos sus puntos de
contacto $c$, con otros granos y con las paredes:
\begin{equation}
  s^{(p)}_{ij} = \frac{1}{A_p} \sum_{c} f^{(c)}_i \, l^{(c)}_j ,
  \qquad
  s^{n(p)}_{ij} = \frac{1}{A_p} \sum_{c} f^{n(c)}_i \, l^{(c)}_j ,
  \label{eq:love}
\end{equation}
donde $\f^{(c)}$ es la fuerza de contacto sobre el grano según
\eqref{eq:fuerza} (con el signo que corresponde a A o a B), $\lv^{(c)}$
es el vector del centro de masa del grano al punto de contacto y
$A_p = \pi R_p^2$ es el \emph{área del propio grano} (no una celda de
Voronoi). $s^n$ usa solo la parte normal de la fuerza,
$\f^{n} = \pm F_n \n$ (con el signo de \eqref{eq:fuerza}). La parte tangencial no se escribe: es $s^t = s - s^n$.

\paragraph{Por qué el área del grano.} Con esta normalización, $s^{(p)}$
es el estrés medio \emph{sobre el grano}. El estrés medio en una región
se obtiene multiplicando por la fracción de área (sección~\ref{sec:phi}).

\paragraph{Relación con el estrés de Cauchy.} Para un cuerpo rígido, el
teorema de la divergencia da
\begin{equation}
  \int_{A_p} \sigma_{ij} \, dA
  = \sum_c f^{(c)}_i l^{(c)}_j
  + \int_{A_p} b_i \, r_j \, dA
  - \int_{A_p} \rho_m \, a_i \, r_j \, dA ,
\end{equation}
con $\vect{b}$ la fuerza de la base por unidad de área, $\vect{r}$ la
posición respecto del centro y $\vect{a}$ la aceleración local. La
fricción de Karnopp es una tracción uniforme y no aporta al primer
momento; el pivoteo y la aceleración angular sí, pero solo a la parte
antisimétrica, y la aceleración centrípeta aporta un término isótropo del
orden de $\omega^2 I$, despreciable. Así, $s^{(p)}$ es el estrés medio de
Cauchy del grano (fórmula de Love) salvo por la parte antisimétrica, que
se analiza en la sección~\ref{sec:antisim}.

\subsection{Formas explícitas para discos}
\label{sec:propiedades}

Con \eqref{eq:fuerza} y \eqref{eq:brazo}, los signos de A y de B se
compensan y se obtiene la misma expresión para los dos cuerpos del
contacto:
\begin{align}
  A_p \, s^{n(p)} &= -\sum_c \rho_c F^{(c)}_n \; \n^{(c)} \otimes \n^{(c)} ,
  \label{eq:sn} \\
  A_p \, s^{t(p)} &= -\sum_c \rho_c F^{(c)}_t \; \tv^{(c)} \otimes \n^{(c)} .
  \label{eq:st}
\end{align}
De aquí salen tres propiedades.

\paragraph{1. La parte normal es simétrica.} $\n \otimes \n$ es simétrico,
así que $s^n_{xy} = s^n_{yx}$.

\paragraph{2. La parte tangencial tiene traza nula.} Como $\tv \cdot \n =
0$, $s^t_{xx} + s^t_{yy} = 0$. La fricción entre granos \emph{no aporta a
la presión}; solo aporta al desviador. La presión del grano es
\begin{equation}
  p^{(p)} = -\tfrac{1}{2}\left(s^{(p)}_{xx} + s^{(p)}_{yy}\right)
  = \frac{1}{2 A_p} \sum_c \rho_c F^{(c)}_n
  \simeq \frac{1}{2 \pi R_p} \sum_c F^{(c)}_n ,
\end{equation}
la suma de las fuerzas normales por unidad de perímetro (a menos de la
corrección $\delta / 2R$ en $\rho$).

\paragraph{3. La parte antisimétrica es el torque de contacto.}
\label{sec:antisim}
Con $\n \times \tv = n_x t_y - n_y t_x = -1$,
\begin{equation}
  s^{(p)}_{yx} - s^{(p)}_{xy} = s^{t(p)}_{yx} - s^{t(p)}_{xy}
  = \frac{1}{A_p} \sum_c \rho_c F^{(c)}_t
  = \frac{\tau^{(p)}_\text{c}}{A_p} ,
  \label{eq:torque}
\end{equation}
donde $\tau_\text{c} = \sum_c (l_x f_y - l_y f_x)$ es el torque de los
contactos respecto del centro. La ecuación de rotación del grano es
$I \dot\omega = \tau_\text{c} + \tau_\text{piv}$. En el promedio temporal
en estado estacionario, $\ave{I\dot\omega} \approx 0$ y
\begin{equation}
  \ave{s^t_{yx} - s^t_{xy}} \approx -\frac{\ave{\tau_\text{piv}}}{A_p} .
\end{equation}
La diferencia entre las curvas de $s^t_{xy}$ y $s^t_{yx}$ mide el torque
que los contactos transmiten a los granos y que la base disipa por
pivoteo. No es un error numérico. Sin fricción con la base, el tensor
sería simétrico en promedio.

La parte simétrica de $s^t$,
$-\frac{1}{2A_p} \sum_c \rho_c F_t (\tv \otimes \n + \n \otimes \tv)$, es
un corte puro con ejes principales a $45^\circ$ de cada normal: es el
aporte de la fricción al corte y a la diferencia de esfuerzos normales.

\paragraph{Verificación.} En la corrida de prueba (\code{frm-04}, una muestra de
34 frames con 13\,600 pares grano-frame), $\max|s^t_{xx} + s^t_{yy}|$ y $\max|s^n_{xy} -
s^n_{yx}|$ son del orden de $10^{-5}$, la precisión con que se escribe el
archivo (seis cifras, con $\max|s| \approx 50$).

\subsection{El archivo \code{.sxy}}

Se escribe cada \code{save\_tensors\_freq} pasos. Cada línea corresponde a
un grano:
\begin{center}
  \code{gID sxx sxy syx syy snxx snxy snyx snyy x y r m vx vy w z\_gg z\_gw}
\end{center}
con posición, radio, masa, velocidad lineal y angular, y los números de
puntos de contacto activos ($F_n > 0$) con granos (\code{z\_gg}) y con
paredes (\code{z\_gw}). La cabecera registra \code{nStep}, \code{n\_frame},
$t$, $\Delta t$, la fase $\omega t \bmod 2\pi$, el commit de git, el
archivo de parámetros y su hash. Con \code{save\_roi\_only: T} se escriben
solo los granos del rectángulo $|x| \le x_\text{roi}$,
$y_\text{min,roi} \le y \le y_\text{max,roi}$; sus tensores incluyen
todos sus contactos, también con granos fuera del rectángulo.

% ---------------------------------------------------------------------------
\section{Perfiles en \texorpdfstring{$y$}{y}: \code{stress\_profile}}
\label{sec:perfiles}

\subsection{Selección y bines}

Se usan los granos con centro en la franja $|x - x_c| \le W$ (en general
$W = D/2 = r$) y $y_\text{min} \le y < y_\text{max}$, en bines de alto
$\Delta y$ (por defecto, un diámetro). Cada frame con tiempo $t$ en
$[t_\text{min}, t_\text{max}]$ se asigna a:
\begin{itemize}
  \item un \emph{bin de fase} $\varphi = \lfloor P \, \theta / 2\pi
    \rfloor$, con $\theta$ la fase de la cabecera y $P$ el número de bines
    de fase (\code{-{}-phase-bins}, por defecto 20);
  \item un \emph{bloque} $b = \lfloor f t \rfloor$, un período de la
    excitación, para los errores.
\end{itemize}
El programa verifica que $\theta$ coincida con $2\pi f t \bmod 2\pi$ (con
la $f$ de \code{-{}-freq}) y que todos los archivos provengan del mismo
commit y del mismo archivo de parámetros.

Sea $\mathcal{B}$ el conjunto de pares grano-frame $(p, \text{frame})$ de
un bin de $y$, y $N_f$ el número de frames usados.

\subsection{Parte de contacto}

Promedio pesado por el área de los granos:
\begin{equation}
  s_{ij} = \frac{\sum_{\mathcal{B}} A_p \, s^{(p)}_{ij}}
                {\sum_{\mathcal{B}} A_p},
  \qquad
  s^n_{ij} = \frac{\sum_{\mathcal{B}} A_p \, s^{n(p)}_{ij}}
                  {\sum_{\mathcal{B}} A_p},
  \qquad
  s^t = s - s^n .
\end{equation}
Con granos de un solo tamaño, es el promedio simple. Por la
sección~\ref{sec:propiedades}, $s^n$ es simétrico, $s^t$ tiene traza nula y
$s_{yx} - s_{xy} = s^t_{yx} - s^t_{xy}$ es el torque medio de contacto por
unidad de área.

\subsection{Parte cinética}
\label{sec:cinetica}

El transporte de momento por las fluctuaciones de velocidad se escribe con
la misma normalización, para que se sume a $s$:
\begin{equation}
  k_{ij} = -\frac{\sum_{\mathcal{B}} m_p \, v'_{p,i} v'_{p,j}}
                 {\sum_{\mathcal{B}} A_p},
  \qquad
  \vv'_p = \vv_p - \vect{u}(y, \varphi),
  \label{eq:k}
\end{equation}
con el signo menos por la convención de compresión negativa. La velocidad
de referencia es la media pesada por masa en el bin de $y$ \emph{y en el
bin de fase}:
\begin{equation}
  \vect{u}(y, \varphi) = \frac{\sum_{\mathcal{B}_\varphi} m_p \vv_p}
                              {\sum_{\mathcal{B}_\varphi} m_p},
\end{equation}
donde $\mathcal{B}_\varphi \subset \mathcal{B}$ son los pares de los
frames con bin de fase $\varphi$. Restar la media por $y$ quita el flujo
medio de descarga; restar la media por fase quita el movimiento coherente
con la oscilación de la base, que de otro modo contaría como agitación.
Con \code{-{}-phase-bins 1}, la referencia es la media temporal.

\paragraph{Cálculo en una pasada.} Por cada celda $(y, \varphi)$ se
acumulan $M = \sum m$, $\vect{Q} = \sum m \vv$ y
$S_{ij} = \sum m v_i v_j$. La suma de las fluctuaciones es
\begin{equation}
  \sum_{\mathcal{B}_\varphi} m \, v'_i v'_j = S_{ij} - \frac{Q_i Q_j}{M},
\end{equation}
y el numerador de \eqref{eq:k} es la suma sobre las $P$ fases
(\code{Acc::fluct}).

$k_{xx}$ y $k_{yy}$ son una temperatura granular traslacional por
componente, en unidades de estrés; su diferencia indica la anisotropía de
la agitación. La rotación no entra en $k$.

\paragraph{Limitaciones.}
\begin{itemize}
  \item \emph{Gradientes dentro de la celda.} La referencia es la media en
    toda la franja y en todo el bin de $y$. Si $\ave{v_y}$ varía con $x$
    dentro de la franja, o con $y$ dentro del bin, esa variación se suma a
    $k$ como falsa agitación. El efecto crece cerca del orificio, donde los
    gradientes son mayores.
  \item \emph{Celdas casi vacías.} Si una celda $(y, \varphi)$ tiene un
    solo grano, $\vv' = 0$ y su aporte es nulo. En los bines con muy pocos
    datos, $k \approx 0$ indica falta de datos, no agitación nula.
\end{itemize}

\subsection{Fracción de área y estrés medio del bin}
\label{sec:phi}

\begin{equation}
  \phi = \frac{\sum_{\mathcal{B}} A_p}{A_\text{bin} \, N_f},
  \qquad A_\text{bin} = 2W \, \Delta y .
\end{equation}
El estrés medio del bin (promedio sobre el área del bin, no sobre los
granos) es
\begin{equation}
  \bar\sigma = \phi \, (s + k) .
\end{equation}
Como cada grano aporta toda su área aunque sobresalga del bin, $\phi$
\emph{no} es una fracción de empaquetamiento física y en franjas angostas
puede superar 1. Es el factor que convierte el estrés sobre los granos en
el estrés medio del bin.

\subsection{Errores por bloques}

Los errores \code{e\_*} suponen que los promedios sobre períodos distintos
son independientes. Para cada bloque $b$ se calcula la media del bloque
con los mismos pesos:
\begin{equation}
  s^{[b]}_{ij} = \frac{\sum_{\mathcal{B} \cap b} A_p \, s^{(p)}_{ij}}
                      {\sum_{\mathcal{B} \cap b} A_p},
\end{equation}
y lo mismo para $s^n$. Con $N_b$ bloques con datos, el error es el error
estándar de la media:
\begin{equation}
  e = \sqrt{\frac{1}{N_b (N_b - 1)} \sum_b \left(s^{[b]} - \bar s\right)^2},
  \qquad \bar s = \frac{1}{N_b} \sum_b s^{[b]} .
\end{equation}

\paragraph{Parte cinética.} Para que las medias por bloque estimen la
misma cantidad que \eqref{eq:k}, las fluctuaciones de cada bloque se miden
respecto de la \emph{misma} referencia global $\vect{u}(y, \varphi)$, no
de una media del bloque (con unos 9 frames por período, esta sería muy
ruidosa). Por cada bloque se guardan $M$, $\vect{Q}$ y $S$ de cada bin de
fase que el bloque toca, y al final
\begin{equation}
  \sum_{\mathcal{B}_\varphi \cap b} m \, v'_i v'_j
  = S^{[b]}_{ij} - Q^{[b]}_i u_j - u_i Q^{[b]}_j + M^{[b]} u_i u_j ,
  \qquad
  k^{[b]}_{ij} = -\frac{\sum_\varphi (\ldots)}{\sum_{\mathcal{B} \cap b} A_p},
\end{equation}
con $\vect{u} = \vect{u}(y, \varphi)$. El error de $k$ se calcula con las
$k^{[b]}$ como el de $s$.

\paragraph{Validez.} Con menos de unos 10 períodos, los errores son poco
confiables. Si hay correlaciones entre períodos consecutivos, los errores
se subestiman; se puede comprobar agrupando bloques más largos.

\subsection{Salidas}

Perfil promediado en fase (\code{-o}), una línea por bin de $y$:
\begin{center}
  \small\code{y n phi vx vy sxx sxy syx syy snxx snxy snyx snyy stxx stxy
  styx styy kxx kxy kyy}\\
  \small\code{e\_sxx e\_sxy e\_syx e\_syy e\_snxx e\_snxy e\_snyx e\_snyy
  e\_kxx e\_kxy e\_kyy}
\end{center}
con \code{n} el número de pares grano-frame y \code{nan} en los bines sin
datos. \code{vx}, \code{vy} son las velocidades medias pesadas por masa.

Perfil resuelto en fase (\code{-{}-phase-output}): las mismas cantidades
(sin $s^t$ ni errores) para cada bin de fase, con la fracción de área
$\phi$ calculada con los frames de esa fase y $k$ con las fluctuaciones de
la celda.

% ---------------------------------------------------------------------------
\section{Gráficos y unidades experimentales}
\label{sec:unidades}

\code{plot\_stress\_profile.py} lee las salidas por el nombre de las
columnas y grafica seis paneles en función de $y$:
\begin{center}
\begin{tabular}{cl}
  \toprule
  Panel & Cantidad \\
  \midrule
  (a) & $p = -(s_{xx} + s_{yy})/2$, error $\frac{1}{2}\sqrt{e_{xx}^2 +
        e_{yy}^2}$ (sin covarianza) \\
  (b) & $-s_{yy}$ \\
  (c) & $-s_{xx}$ \\
  (d) & $s_{xy}$ (línea), $s_{yx}$ (trazos) \\
  (e) & $s^t_{xy}$ (línea), $s^t_{yx}$ (trazos) \\
  (f) & $-k_{yy}$ (línea), $-k_{xx}$ (trazos), con sus errores \\
  \bottomrule
\end{tabular}
\end{center}
Con \code{-{}-continuo}, el estrés y sus errores se multiplican por $\phi$
(estrés medio del bin).

\paragraph{Conversión} (\code{unidades.py}, valores de
\code{utils/reduced\_units.ods}):
\begin{center}
\begin{tabular}{lll}
  \toprule
  Magnitud & Unidad reducida & Valor \\
  \midrule
  longitud & $d$ & $0.005$ m \\
  masa & $m$ & $2.10 \times 10^{-4}$ kg \\
  tiempo & $\sqrt{d/g}$ & $0.0225877$ s \\
  velocidad & $\sqrt{gd}$ & $0.221359$ m/s \\
  fuerza & $mg$ & $0.002058$ N \\
  estrés 2D & $mg/d$ & $0.4116$ N/m \\
  \bottomrule
\end{tabular}
\end{center}

\paragraph{Inversión del eje $y$.} Al pasar a las coordenadas del
experimento, $y \to -y$. Cambian de signo las cantidades con un único
índice $y$: $v_y$ y las componentes cruzadas $xy$, $yx$ de $s$, $s^n$,
$s^t$ y $k$. Las diagonales y los errores no cambian.

% ---------------------------------------------------------------------------
\section{Verificaciones}

\begin{itemize}
  \item Balance de impulso por grano (\code{check\_balance\_freq}):
    residuo relativo $\sim 10^{-5}$ lineal y $\sim 10^{-6}$ angular sin
    TOI; valida la convención de la tangente.
  \item Identidades de discos (sección~\ref{sec:propiedades}): se cumplen a la
    precisión de la salida.
  \item \code{stress\_profile} contra una implementación independiente en
    Python: diferencias relativas $\le 3 \times 10^{-7}$ en $s$, $s^n$,
    $k$ y en los errores de $k$.
\end{itemize}

\end{document}
